\documentclass[aspectratio=169, 11pt]{beamer}
\usepackage{fontspec}
\setmainfont{TeX Gyre Termes}
\setsansfont{TeX Gyre Heros}
\setmonofont{Latin Modern Mono}
\usepackage[spanish,es-nodecimaldot]{babel}
\usepackage{amsmath, amssymb, amsfonts, booktabs, graphicx}
\usepackage{tikz}
\usetikzlibrary{shapes.geometric, arrows.meta, positioning, calc, fit, shadows}
\usetheme{Madrid}
\usecolortheme{beaver}
\definecolor{unalRojo}{RGB}{148, 36, 36}
\definecolor{verdeOK}{RGB}{34, 139, 34}
\definecolor{naranjaAlerta}{RGB}{230, 126, 34}
\definecolor{rojoAlerta}{RGB}{192, 57, 43}
\setbeamercolor{structure}{fg=unalRojo}
\setbeamercolor{block title}{bg=unalRojo!15, fg=unalRojo}
\setbeamercolor{block body}{bg=unalRojo!5}
\setbeamercolor{alerted text}{fg=rojoAlerta}
\newcommand{\BSF}{\textit{Hermetia illucens}}
\newcommand{\MAPE}{\text{MAPE}}
\newcommand{\DEB}{\text{DEB}}
\newcommand{\NDIR}{\text{NDIR}}
\newcommand{\MRV}{\text{MRV}}
\setbeamertemplate{navigation symbols}{}
\setbeamertemplate{footline}[frame number]

\begin{document}

% ---------------------------------------------------------------------------
% Slide puente: De lo mecanic\'{i}stico a lo h\'{i}brido
% ---------------------------------------------------------------------------
\begin{frame}{El mecanic\'{i}stico asume un mundo homog\'{e}neo que no existe}
    \begin{columns}[T]
        \begin{column}{0.55\textwidth}
            \begin{alertblock}{Limitaciones estructurales}
                \begin{itemize}
                    \item Las larvas se agregan como fluidos activos: gradientes 3D de T y O$_2$.
                    \item No discrimina CO$_2$ de larva vs. microbiota del sustrato.
                    \item Ceguera total ante anaerobiosis: no predice CH$_4$.
                    \item Validaci\'{o}n destructiva altera el reactor mismo.
                \end{itemize}
            \end{alertblock}
        \end{column}
        \begin{column}{0.4\textwidth}
            \begin{center}
                \begin{tikzpicture}[scale=0.7, transform shape]
                    \draw[thick, fill=gray!10] (0,0) rectangle (3.5,3);
                    \foreach \x/\y in {0.4/0.4, 0.9/0.2, 0.6/0.9, 2.8/2.5, 2.2/2.8, 3/2, 1.5/2.2, 2.5/1.2, 0.2/2.5, 1.2/1.5} {
                        \fill[unalRojo] (\x,\y) circle (0.12);
                    }
                    \draw[thick, dashed, rojoAlerta, fill=rojoAlerta!15] (2.6,1.8) circle (0.7);
                    \node[font=\footnotesize, rojoAlerta] at (2.6,1.8) {An\'{o}xico};
                \end{tikzpicture}
            \end{center}
        \end{column}
    \end{columns}

    \vfill
    \begin{center}
        \textbf{Para cerrar la brecha sin romper la f\'{i}sica $\rightarrow$ arquitectura h\'{i}brida}
    \end{center}
\end{frame}

% ---------------------------------------------------------------------------
% Slide 2: Sensor virtual
% ---------------------------------------------------------------------------
\begin{frame}{Sensor virtual de biomasa}
    \begin{center}
        \begin{tikzpicture}[scale=0.85, transform shape,
            input/.style={rectangle, draw=unalRojo, fill=unalRojo!5, minimum width=2.6cm, minimum height=0.85cm, align=center, font=\small},
            nn/.style={rectangle, draw=gray, fill=gray!15, minimum width=2.6cm, minimum height=1.8cm, align=center, font=\small},
            output/.style={rectangle, draw=verdeOK, fill=verdeOK!10, minimum width=2.6cm, minimum height=0.85cm, align=center, font=\small}]

            \node[input] (d) at (0,1.3) {Densidad inicial};
            \node[input] (t) at (0,0) {Temperatura};
            \node[input] (h) at (0,-1.3) {Humedad};

            \node[nn] (red) at (4.5,0) {Red neuronal densa};

            \node[output] (b) at (9,0) {$B(t)$ estimada};

            \draw[-{Stealth}, thick] (d) -- (red);
            \draw[-{Stealth}, thick] (t) -- (red);
            \draw[-{Stealth}, thick] (h) -- (red);
            \draw[-{Stealth}, thick] (red) -- (b);
        \end{tikzpicture}
    \end{center}

    \vfill
    \begin{itemize}
        \item Resuelve el problema inverso: infiere el estado interno \textbf{sin muestreos destructivos}.
        \item Actualiza en cada paso el t\'{e}rmino $dB/dt$ de la EDO de Eriksen.
    \end{itemize}
\end{frame}

% ---------------------------------------------------------------------------
% Slide 3: PINN
% ---------------------------------------------------------------------------
\begin{frame}{PINN: correcci\'{o}n f\'{i}sicamente consistente}
    \begin{center}
        \begin{tikzpicture}[scale=0.8, transform shape,
            box/.style={rectangle, rounded corners, draw=unalRojo, fill=unalRojo!5, minimum width=3.2cm, minimum height=1.1cm, align=center, font=\small},
            constraint/.style={rectangle, rounded corners, draw=verdeOK, fill=verdeOK!10, minimum width=3.2cm, minimum height=1.1cm, align=center, font=\small}]

            \node[box] (data) at (0,0) {Datos sensores CO$_2$, CH$_4$};
            \node[box] (pinn) at (5.5,0) {PINN};
            \node[constraint] (edo) at (5.5,-2.5) {EDO Eriksen (restricci\'{o}n dura)};
            \node[box] (out) at (11,0) {Predicci\'{o}n ajustada};

            \draw[-{Stealth}, thick] (data) -- (pinn);
            \draw[-{Stealth}, thick] (pinn) -- (out);
            \draw[-{Stealth}, thick, dashed, verdeOK] (edo) -- (pinn);
        \end{tikzpicture}
    \end{center}

    \vfill
    \begin{block}{Funci\'{o}n de p\'{e}rdida}
        \centering
        $\displaystyle \mathcal{L} = \underbrace{\mathcal{L}_{\text{datos}}}_{\text{Error vs. sensores}} + \underbrace{\mathcal{L}_{\text{EDO}}}_{\text{Residual f\'{i}sico}}$
    \end{block}
\end{frame}

% ---------------------------------------------------------------------------
% Slide 4: Justificaci\'{o}n CO2/CH4
% ---------------------------------------------------------------------------
\begin{frame}{CO$_2$ y CH$_4$: se\~{n}ales del metabolismo y del fallo}
    \begin{columns}[T]
        \begin{column}{0.48\textwidth}
            \begin{block}{CO$_2$ como se\~{n}al fisiol\'{o}gica}
                \begin{itemize}
                    \item Gas primario de la respiraci\'{o}n aerobia.
                    \item Vinculado a $B(t)$ por el coeficiente $Y_{\text{CO}_2}$.
                    \item Variable de estado para calibrar el h\'{i}brido.
                \end{itemize}
            \end{block}
        \end{column}
        \begin{column}{0.48\textwidth}
            \begin{alertblock}{CH$_4$ como se\~{n}al de fallo}
                \begin{itemize}
                    \item Indica anaerobiosis en microambientes.
                    \item Detecci\'{o}n temprana esencial para la gesti\'{o}n.
                    \item PGW 27,2$\times$ CO$_2$ (IPCC AR6).
                \end{itemize}
            \end{alertblock}
        \end{column}
    \end{columns}

    \vfill
    \begin{center}
        \textbf{N$_2$O queda fuera del alcance cuantitativo:} requiere sensores electroqu\'{i}micos espec\'{i}ficos y cadenas de calibraci\'{o}n independientes, fuera del dominio de validaci\'{o}n de esta fase.
    \end{center}
\end{frame}

\end{document}
